\section{Trace-preserving polynomial flows}
\label{sec:flow}

Voiculescu's exponentiation theorem turns trace-preserving polynomial
vector fields on algebraically free self-adjoint generators into
one-parameter automorphism groups~\cite[Theorem~3.8]{voiculescu2002cyclomorphy}.
We prove the version needed here directly for Haar unitaries and a fixed
bounded logarithm. The proof treats the unitary relations explicitly
and derives preservation of all evaluated relations from moment invariance.

Retain the notation of Section~\ref{sec:prelim}. Thus
$(A_1,\ldots,A_n,C)$ is a freely generating Haar tuple in the finite
tracial von Neumann algebra $(M,\tau)$,
$\Gamma=\F_n=\langle x_1,\ldots,x_n\rangle$, and
\[
 S=\arg(C),\qquad S_g=A_gSA_g^*,\qquad
 s(k)=\sum_{g\in\Gamma}k(g)S_g.
\]
Write $\R^{(\Gamma)}$ for the finitely supported real functions on
$\Gamma$, and let $\e\in\ell^2(\Gamma)$ be the unit vector at the
identity. All coefficient vectors used in this section are finitely supported.
The left regular action on them is
$(\lambda(g)k)(b)=k(g^{-1}b)$.

Choose real coefficient vectors $h_1,\ldots,h_n$. There is a unique
cocycle $D:\Gamma\to\R^{(\Gamma)}$ satisfying
\begin{equation}
 D_{x_j}=h_j,\qquad D_{gb}=D_g+\lambda(g)D_b.
 \label{eq:flow-cocycle}
\end{equation}
Indeed, set $D_{x_j^{-1}}=-\lambda(x_j^{-1})h_j$ and define the value
on a word by summing the shifted values of its letters. The two terms
associated with an adjacent inverse pair cancel, so the definition
respects reduction. The cocycle identity and uniqueness follow.

The next proposition turns small $\ell^2$ norms of the $h_j$ into
small operator-norm displacements of the generators. For a word
$w\in\Gamma$, a small $\|D_w-\e\|_{\ell^2}$ also makes the
time-one image of $A_w$ close to $CA_w$.

\begin{proposition}[Polynomial flow and displacement bounds]
\label{prop:flow}
There is a one-parameter group $(\beta_t)_{t\in\R}$ of unital
trace-preserving normal $*$-automorphisms of $M$ that fixes $S$ and
whose restriction to
\[
 B=C^*(A_1,\ldots,A_n,S)
\]
is point-norm continuous. Writing $A_j(t)=\beta_t(A_j)$ and evaluating
group words at this tuple, the paths satisfy
\begin{equation}
 \frac{d}{dt}A_j(t)=\ii s_t(h_j)A_j(t),\qquad
 s_t(k)=\sum_{g\in\Gamma}k(g)A_g(t)SA_g(t)^*,\qquad A_j(0)=A_j.
 \label{eq:flow-ode}
\end{equation}
The derivatives here are in operator norm. With $K=3\pi$ as in
\eqref{eq:free-sum}, for every $t\in\R$,
\begin{align}
 \norm{\beta_t(A_j)-A_j}
 &\leq |t|K\|h_j\|_{\ell^2},
 \label{eq:generator-displacement}\\
 \norm{\beta_t(A_w)-\exp(\ii tS)A_w}
 &\leq |t|K\|D_w-\e\|_{\ell^2}
 \qquad(w\in\Gamma).
 \label{eq:word-displacement}
\end{align}
In particular, at $t=1$ the target in
\eqref{eq:word-displacement} is $CA_w$.
\end{proposition}

We first establish a trace identity on a formal polynomial algebra.
This will prove invariance along the actual flow without assuming
freeness of the evolving tuple.

\subsection{The formal derivation and its trace identity}

Let $\mathcal P$ be the complex unital $*$-algebra generated by
unitaries $a_1,\ldots,a_n$ and a self-adjoint symbol $z$, subject only
to these unitary and self-adjoint relations. Put
\[
 z_g=a_gza_g^*,\qquad \sigma(k)=\sum_gk(g)z_g.
\]
Here $a_g$ is the group word in the formal unitaries. Evaluation at
$(A_1,\ldots,A_n,S)$ defines a $*$-homomorphism
$\operatorname{ev}_0:\mathcal P\to B$; we make no assumption about
its kernel.

Define a complex-linear derivation by the rules
\begin{equation}
 \delta z=0,\qquad
 \delta a_j=\ii\sigma(h_j)a_j,\qquad
 \delta a_j^*=-\ii a_j^*\sigma(h_j).
 \label{eq:formal-derivation}
\end{equation}
These rules define a $*$-derivation on $\mathcal P$. The reality of
$h_j$ makes $\sigma(h_j)$ self-adjoint, so the rules respect
adjoints. They also respect the defining unitary relations: for
example,
\[
 \delta(a_j^*a_j)
 =-\ii a_j^*\sigma(h_j)a_j+\ii a_j^*\sigma(h_j)a_j=0,
\]
and the calculation for $a_ja_j^*$ is the same.
The identity
$a_g\sigma(k)a_g^*=\sigma(\lambda(g)k)$ and the cocycle rule give
\begin{equation}
 \delta a_g=\ii\sigma(D_g)a_g,
 \qquad
 \delta z_g=\ii[\sigma(D_g),z_g].
 \label{eq:formal-word-derivative}
\end{equation}
For instance, multiplying the formulas for $g$ and $b$ gives the
formula for $gb$ with coefficient $D_g+\lambda(g)D_b$; inverse
letters are covered by the unitary relations.

\begin{lemma}[Universal initial trace identity]
\label{lem:flow-trace-identity}
For every $p\in\mathcal P$ and every integer $r\geq1$,
\begin{equation}
 \tau\bigl(\operatorname{ev}_0(\delta^r p)\bigr)=0.
 \label{eq:universal-trace-identity}
\end{equation}
\end{lemma}

\begin{proof}
First take $r=1$. Give a monomial its group label by ignoring each
occurrence of $z$ and multiplying the remaining group symbols in
order. Each term of $\sigma(k)$ has label $e$. The replacement
rules~\eqref{eq:formal-derivation} consequently preserve a monomial's
label term by term. If that label is nonidentity,
Lemma~\ref{lem:labels} makes every evaluated derivative term have
trace zero.

An identity-label monomial is a polynomial in the $z_g$. To see
this algebraically, repeatedly use $a_gz=z_ga_g$ to move group
factors to the right; the final group factor is the identity.
Because these are identities in $\mathcal P$ and $\delta$ is a
derivation on that algebra, differentiating this rewriting gives
the same result as differentiating the original monomial.
It is therefore enough to consider
$p=F(z_g:g\in I)$, with $I\subset\Gamma$ finite and $F$ a
noncommutative polynomial.

For a fixed $g\in I$, remove the commuting own-variable term from
the commutator in~\eqref{eq:formal-word-derivative}. Its evaluated
direction is
\[
 \operatorname{ev}_0(\delta z_g)=\ii[Y_g,S_g],\qquad
 Y_g=\sum_{b\ne g}D_g(b)S_b.
\]
The element $Y_g$ is a bounded self-adjoint in
$B_g=W^*(S_b:b\ne g)$, and $W^*(S_g)$ is free from $B_g$ by
Lemma~\ref{lem:conjugates}.

We use the following direct consequence of freeness. If a unital
subalgebra $\mathcal A$ is free from a unital subalgebra
$\mathcal B$ and $U\in\mathcal B$ is unitary, then
$U\mathcal A U^*$ is free from $\mathcal B$. Indeed, conjugating
an alternating product of centered elements from these two
algebras by $U^*$ gives an alternating product of centered
elements from $\mathcal A$ and $\mathcal B$: a fixed centered
$b\in\mathcal B$ becomes $U^*bU$, which is still centered in
$\mathcal B$. Its trace is zero. Conjugation preserves the
marginal law of each element of $\mathcal A$. Since freeness
determines mixed moments from the marginal laws, it follows that
replacing $S_g$ by $US_gU^*$ preserves its joint law with any
fixed collection of elements of $B_g$.

Apply this observation with $U=\exp(\ii uY_g)$ for real $u$.
The scalar function obtained by replacing every occurrence of
$S_g$ in $F$ by $\exp(\ii uY_g)S_g\exp(-\ii uY_g)$ and leaving
all other variables fixed has constant trace. Its derivative at
zero is the trace of the sum over those occurrences with the
direction $\ii[Y_g,S_g]$ inserted. Sum these zero derivatives over
$g\in I$. The product rule and
\eqref{eq:formal-word-derivative} identify the sum exactly with
$\tau(\operatorname{ev}_0(\delta p))$, proving the case $r=1$.
Linearity handles arbitrary finite sums of monomials.

The result just proved is the identity
$\tau\circ\operatorname{ev}_0\circ\delta=0$ on all of
$\mathcal P$. Since $\delta$ maps finite polynomials to finite
polynomials, applying it to $\delta^{r-1}p$ proves
\eqref{eq:universal-trace-identity}. In this application every
derivative of every coefficient polynomial is included. The
auxiliary conjugations above are used only to compute a first
variation at the initial tuple; they need neither lift to motions
of the $a_j$ nor be jointly integrated.
\end{proof}

\subsection{Global solutions and actual moment invariance}

We now solve the polynomial equation globally in $B$ and use
analyticity to turn the initial trace identity into moment invariance.
For a tuple $X=(X_1,\ldots,X_n)$, let $X_g$ denote the reduced
group word evaluated in a $*$-algebra using adjoints for inverse
letters. This notation defines a polynomial even for nonunitary
tuples, so the right side of~\eqref{eq:flow-ode} defines the finite
polynomial vector field
\[
 \mathcal V_j(X)=\ii H_j(X)X_j,
 \qquad
 H_j(X)=\sum_g h_j(g)X_gSX_g^*.
\]

\begin{lemma}[Global solvability in a generated subalgebra]
\label{lem:flow-global-ode}
Let $E$ be any closed unital $C^*$-subalgebra of $M$ containing
$S$, and let $V_1,\ldots,V_n$ be unitaries in $E$. The equation
$X'=\mathcal V(X)$ with $X(0)=V$ has a unique global solution
in $E^n$ that is smooth in norm, and all its coordinates remain unitary.
\end{lemma}

\begin{proof}
Regard $E^n$, with its maximum norm, as a real Banach space.
Adjunction is bounded and real-linear, so $\mathcal V$ is a
finite smooth polynomial map. Its norm and its Lipschitz constant
are bounded on every fixed bounded ball. Picard iteration on a
sufficiently short interval gives a unique local solution;
repeated differentiation makes it smooth.

Each $H_j(X)$ is self-adjoint, even before unitarity is known.
Along a local solution,
\[
 \frac{d}{dt}(X_j^*X_j)
 =-\ii X_j^*H_j(X)X_j+\ii X_j^*H_j(X)X_j=0.
\]
Thus $X_j^*X_j=1$. Such an isometry is unitary in $M$: the
projection $1-X_jX_j^*$ has trace
$1-\tau(X_j^*X_j)=0$, and faithfulness makes it zero.
This equality also holds in $E$, so $\norm{X_j}=1$ throughout
the local solution.

On this unitary trajectory,
\[
 \norm{X_j'}\leq \norm{S}\sum_g|h_j(g)|.
\]
If a maximal interval had a finite endpoint, this bounded
derivative would make the tuple norm-Cauchy as time approached
that endpoint. Its limit lies in the closed algebra $E$ and is
again a unitary tuple. Local existence at the limit continues
the solution, a contradiction. The same argument applies at
either endpoint, giving a global solution. Uniqueness follows
from local uniqueness on overlapping intervals. This
continuation uses completeness and bounded derivative, not
compactness of a bounded ball.
\end{proof}

Apply Lemma~\ref{lem:flow-global-ode} in $E=B$ with initial
tuple $A$. Let $A(t)$ be its solution and let
$\operatorname{ev}_t:\mathcal P\to B$ evaluate the formal
generators at $A(t),S$. Unitarity makes this a well-defined
evaluation on the formal algebra for each $t$.

\begin{lemma}[Moment invariance]
\label{lem:flow-moments}
For every $p\in\mathcal P$ and $t\in\R$,
\begin{equation}
 \tau\bigl(\operatorname{ev}_t(p)\bigr)
 =\tau\bigl(\operatorname{ev}_0(p)\bigr).
 \label{eq:flow-moments}
\end{equation}
\end{lemma}

\begin{proof}
The differential equation and the product rule give, for every
integer $r\geq0$,
\begin{equation}
 \frac{d^r}{dt^r}\operatorname{ev}_t(p)
 =\operatorname{ev}_t(\delta^r p).
 \label{eq:flow-iterated-derivatives}
\end{equation}
We verify that these derivatives give an actual convergent
Taylor expansion for the trace.

Choose fixed polynomial representatives for the derivatives
of the formal letters in~\eqref{eq:formal-derivation}. There
are constants $J\geq1$ and $B_1\geq1$ such that replacing any
one letter by its derivative increases monomial length by at
most $J$ and produces total absolute coefficient sum at most
$B_1$. These constants are finite because the $h_j$ have finite
support. If a representative of $p$ has degree at most $l$, where
$l\geq1$, and total absolute coefficient sum $b_0$, repeated
replacement gives a representative of $\delta^r p$ with degree
at most $l+rJ$ and coefficient sum at most
\[
 b_0 B_1^r\prod_{v=0}^{r-1}(l+vJ).
\]
Let $R=\max(1,\norm{S})$. All the evaluated group letters are
unitary, so~\eqref{eq:flow-iterated-derivatives} implies, uniformly
for real $t$,
\begin{align}
 \frac{1}{r!}
 \left\|\frac{d^r}{dt^r}\operatorname{ev}_t(p)\right\|
 &\leq b_0R^l(B_1R^J)^r
       \frac{\prod_{v=0}^{r-1}(l+vJ)}{r!} \notag\\
 &\leq b_0R^l\bigl[B_1R^J(l+J)\bigr]^r.
 \label{eq:flow-analytic-bound}
\end{align}
The last inequality uses $l+vJ\leq(l+J)(v+1)$.

Set $f_p(t)=\tau(\operatorname{ev}_t(p))$ and
$c_p=B_1R^J(l+J)$. The integral form of Taylor's remainder and
\eqref{eq:flow-analytic-bound} bound the remainder after degree
$r$ at increment $u$ by $b_0R^l(c_p|u|)^{r+1}$. Consequently
$f_p$ is real-analytic near every real time, with positive
radius at least $c_p^{-1}$. By
Lemma~\ref{lem:flow-trace-identity}, all its positive-order
derivatives at zero vanish. It is constant near zero and hence,
by real-analytic continuation on the connected line $\R$,
constant everywhere.

The constants here may depend on the fixed field and on $p$.
No radius common to all polynomials or to different choices of
the $h_j$ is required. The bounds used only unitarity of the
solution and boundedness of $S$; freeness at positive times has
not been assumed.
\end{proof}

\subsection{Surjectivity, normal extension, and estimates}

\begin{proof}[Proof of Proposition~\ref{prop:flow}]
For a fixed real $t$, polynomial evaluation defines a map
\[
 \varphi_t:\operatorname{ev}_0(\mathcal P)
      \longrightarrow B,
 \qquad
 \varphi_t\bigl(\operatorname{ev}_0(p)\bigr)
      =\operatorname{ev}_t(p).
\]
This is well defined. If $\operatorname{ev}_0(p)=0$, then
\eqref{eq:flow-moments} applied to $p^*p$ gives
\[
 \tau\bigl(\operatorname{ev}_t(p)^*
                 \operatorname{ev}_t(p)\bigr)=0,
\]
so faithfulness implies $\operatorname{ev}_t(p)=0$.
Thus any additional relations at the evaluated initial tuple
are respected as a consequence of actual moment invariance.
The map is unital, complex-linear, multiplicative, and
adjoint-preserving because formal evaluation has those
properties.

For an element $b$ of a finite von Neumann algebra with faithful
tracial state,
\[
 \norm{b}
 =\lim_{k\to\infty}\tau((b^*b)^k)^{1/(2k)}.
\]
For completeness, the upper bound is immediate; for the lower
bound, every nonzero spectral projection of $b^*b$ sufficiently
near its norm has positive trace by faithfulness. Applying
\eqref{eq:flow-moments} to $(p^*p)^k$ therefore shows that
$\varphi_t$ is isometric. It extends uniquely to a
trace-preserving isometric $*$-homomorphism on $B$, with range
\[
 B_t=C^*(A_1(t),\ldots,A_n(t),S).
\]
Indeed, its range is closed and is the norm closure of the
polynomials in these time-$t$ generators.

We next prove $B_t=B$, which establishes surjectivity.
Apply Lemma~\ref{lem:flow-global-ode} inside $B_t$ to the
initial tuple $A(t)$, using the same fixed element $S$ and the
same polynomial field. Its global solution, viewed in $M$,
must be $u\mapsto A(t+u)$ by uniqueness. At $u=-t$ it recovers
the original tuple. Thus every $A_j$ belongs to $B_t$, and
$S$ already belongs to $B_t$, giving $B\subseteq B_t$.
The forward construction gives the reverse inclusion. Hence
$\varphi_t$ is a $C^*$-automorphism of $B$ fixing $S$.

Since $\exp(\ii S)=C$, the von Neumann algebra generated by
$B$ is $M$. Its vectors are therefore dense in $L^2(M,\tau)$,
as in the density argument in
Lemma~\ref{lem:haar-identification}. Trace preservation and
surjectivity define a Hilbert-space unitary
\[
 U_t[b]=[\varphi_t(b)]\qquad(b\in B)
\]
on $L^2(M,\tau)$, fixing $[1]$. If $L_b$ denotes left
multiplication, then on the dense set of vectors coming from
$B$,
\[
 U_t L_b U_t^*=L_{\varphi_t(b)}\qquad(b\in B).
\]
The faithful normal left representation identifies the von
Neumann closure of $L(B)$ with $M$. Conjugation by $U_t$
therefore gives a normal $*$-automorphism $\beta_t$ of $M$
extending $\varphi_t$. It is trace-preserving because $U_t$
fixes $[1]$, and it fixes $S$. This normal extension is unique
by ultraweak density of $B$.

The maps form a one-parameter group. For fixed $s$, the path
$t\mapsto\varphi_s(A(t))$ satisfies the same polynomial
equation, because $\varphi_s$ is a bounded $*$-homomorphism
fixing $S$. Its initial tuple is $A(s)$, so uniqueness gives
\[
 \varphi_s(A_j(t))=A_j(s+t).
\]
The equality on the generators and on $S$ yields
$\varphi_s\varphi_t=\varphi_{s+t}$ on $B$, and uniqueness
of normal extension gives $\beta_s\beta_t=\beta_{s+t}$ on
$M$. Also $\beta_0$ is the identity. Norm continuity on
polynomials in the generators follows from the solution
paths; density and the common isometry bound extend it to
point-norm continuity on $B$. No point-norm continuity on
arbitrary elements of $M$ is needed.

It remains to prove the bounds. For each finitely supported
real $k$,
\[
 s_t(k)=\beta_t(s(k)),\qquad
 \norm{s_t(k)}=\norm{s(k)}\leq K\|k\|_{\ell^2}
\]
by~\eqref{eq:free-sum}. Integrating
$\norm{A_j'(t)}\leq K\|h_j\|_{\ell^2}$ between zero and
$t$ proves~\eqref{eq:generator-displacement}.

For a fixed $w\in\Gamma$, evaluation of
\eqref{eq:formal-word-derivative} along the flow gives
\[
 \frac{d}{dt}A_w(t)=\ii s_t(D_w)A_w(t).
\]
Moreover $s_t(\e)=S$. Direct differentiation, with the order
of factors retained, yields
\[
 \frac{d}{dt}\bigl(\exp(-\ii tS)A_w(t)\bigr)
 =\ii\exp(-\ii tS)\bigl(s_t(D_w)-S\bigr)A_w(t).
\]
The exterior factors are unitary, so its norm is at most
$K\|D_w-\e\|_{\ell^2}$. Integrating between zero and $t$
and multiplying on the left by $\exp(\ii tS)$ proves
\eqref{eq:word-displacement}.
\end{proof}
